\section{Related Work}
\label{sec:related}

\paragraph{Compressed LLMs as agents.} ACBench \citep{dong2025acbench} evaluates 15 compressed models at temperature 0 (tool use drops 1--3\% at 4 bits, real-world applications 10--15\%); our notes record one run per configuration and no reported variance. The closest work is \emph{Flat Score, Amplified Failures} \citep{jang2026flat}: BF16, FP8 and AWQ-INT4 agents on \tautwo{} retail and telecom with four trials per task. It finds no score change surviving multiple-comparison correction, reports TOST equivalence at $\pm 7.5$ points with a noise floor (about 1.5 points) from one repeat run, and shows failure modes rising up to $2.5\times$ while the aggregate stays flat; a two-error budget re-exposes a 17-point gap. Its margin account of flipped decisions is verbal (token log-probabilities are not in its logs), and it measures neither divergence nor a perturbation null; we add both, plus measured margins (Table~\ref{tab:comparison}, Appendix~\ref{app:comparison}). Other work reports token inflation under low-bit quantization \citep{lian2026quantization}, task-dependent precision needs of quantized agents \citep{lu2026mixquant,zhang2026quantclaw}, function calling robust to aggressive quantization \citep{taheri2026iffc}, and closed-loop evaluation of quantized vision-language-action policies \citep{xu2026vlaquantbench}.

\paragraph{Evaluation noise and reference nulls.} See \citet{heineman2025signal,kapoor2024agentsmatter,yao2024tau} for benchmark separability, agent benchmarking practice and pass\^{}$k$ reliability. On \tautwo{}, an analysis of released trajectories estimates that single runs must differ by 7--8 (retail) or 10--11 (airline) points to exceed noise under sampling \citep{tau2issue540}; single-run pass@1 of SWE-bench Verified agents varies by 2.2--6.0 points, with SDs above 1.5 points even at temperature 0 \citep{bjarnason2026randomness}. On single-turn BFCL, prompt perturbations produce paired SDs 11--58$\times$ those of reruns, with multi-turn agents out of scope \citep{chen2026noisefloor}. \citet{singh2026certifying} finds per-item churn about $5\times$ the net delta for compressed models on static benchmarks and proposes paired equivalence testing at a \emph{declared} margin ($\pm 2$ points in its example). The Replay Gap \citep{gonuguntla2026replay} uses same-model control forks as a null for model switching in SWE agents; FP8-served forks diverge in over 90\% of cases while AWQ-served ones stay near-identical, so rerun-null width depends on the serving format. Statistical precedents include composite perturbation nulls \citep{acharyya2025composite}, data-dependent equivalence margins \citep{koobs2026equivalence} and counterfactual-prompting baselines \citep{yang2026comparedtowhat}. Our null comes from the reference model's \emph{own} perturbations, task by task, in a closed loop, with its width explained by a margin model.

\paragraph{The per-step margin model (inherited).} \citet{du2026greedy} show that BF16 and FP16 greedy decoding diverge on 49--100\% of prompts and that a token flips when the perturbation of the logit gap exceeds the top-two margin, with the heuristic $P(\text{flip at } t) \lesssim P(\Delta_t < \sigma_z)$ (their Eq.~14). The same relation holds for INT8 kernel swaps (margin-to-flip ROC-AUC 0.942; \citealp{chen2026integer}), and low margins can be exploited to restore determinism under batching \citep{chu2026margingate}; numerical nondeterminism in serving is documented by \citet{yuan2025nondeterminism} and \citet{he2025defeating}. We take the per-step model as given; our contribution is its composition over a closed loop, the resulting floor and two-regime structure along one perturbation axis, the later damage threshold, and its use to define and interpret a null.

\paragraph{Divergence in multi-turn settings; open vs.\ closed loop.} Prefix caching changes 36.2\% of multi-turn tool episodes at 16 bits and 75.0\% at 4 bits, with saturation at the coarsest settings and a 3.5\% single-turn-bridge rate explained verbally by compounding \citep{patodiya2026same}. Meaning-bearing rewrites destabilize agents more than presentation changes \citep{zhang2026agentdiff} (our nonce is a presentation change); multi-turn loss is mostly unreliability rather than aptitude \citep{laban2025lost}; long-horizon execution compounds errors \citep{sinha2026illusion}; and classical work on compounding error in sequential prediction \citep{ross2011dagger,bengio2015scheduled} motivates the closed-loop view. Equal accuracy can hide changed predictions \citep{dutta2024accuracy,deiseroth2023divergent}, distribution-divergence metrics complement accuracy \citep{qamar2026divergence}, and in driving, open-loop metrics correlate with closed-loop scores only non-monotonically \citep{wang2026openloopdriving}. We add a direct open- versus closed-loop measurement for quantized tool agents (Appendix~\ref{app:deepdive}).
